\chapter{THE TRANSCENDENCY OF $a^b$ FOR IRRATIONAL ALGEBRAIC $b$ AND ALGEBRAIC $a \neq 0, 1$}

Let $p$ be any complex number $\neq 0$ such that $e^p = a$ is algebraic, and let $\beta$ be any irrational number such that also $e^{p\beta} = a^\beta = c$ is algebraic. Because of the addition theorem $f(x+y) = f(x)f(y)$ for the function $f(x) = e^{px} = a^x$ it follows that $f(x)$ takes the algebraic values $a^p c^q$ for all $x = p + \beta q$ ($p, q = 0, \pm 1, \pm 2, \dots$). These $x$ are different from each other, because of the irrationality of $\beta$.

In 1929 Gelfond made the important discovery that $\beta$ cannot be a non-real quadratic irrationality; this means that the number $a^\beta = e^{\beta \log a}$ is transcendental for every imaginary quadratic irrationality $\beta$ and every algebraic $a \neq 0, 1$, except for $\log a = 0$. In his proof Gelfond applied the interpolation formula of \S 14, Chapter I, to the special case $f(z) = a^z = e^{z \log a}$ and the sequence $z_1, z_2, \dots$ consisting of the points $p + \beta q$ ($p, q = 0, \pm 1, \pm 2, \dots$) in a certain arrangement. The coefficients $a_0, a_1, \dots$ in the expansion
\[
f(z) = a_{0}F_{0}(z) + a_{1}F_{1}(z) + \dots,
\]
\[
F_{n}(z) = \prod_{l=1}^{n} (z-z_{l})
\]
are given by
\[
a_{n-1} = \frac{1}{2\pi i} \int_{C} \frac{f(z)}{F_{n}(z)} \, dz = \sum_{k=1}^{n} \frac{f(z_{k})}{F_{n}'(z_{k})},
\]
\[
F_{n}'(z_{k}) = \prod_{\substack{l=1\\l \neq k}}^{n} (z_{k} - z_{l}),
\]
and this shows that they are numbers in the field $K$ generated by $a, \beta, a^\beta$. On one hand it follows from the integral formula for $a_{n-1}$ that the coefficients tend rapidly to $0$ as $n \to \infty$; on the other hand, if $K$ were an algebraic number field, one can obtain an estimate of $|a_n|$ and the denominator of $a_n$. In case of an imaginary quadratic $\beta$ these estimates imply that $a_n = 0$ for all sufficiently large $n$. This is contradictory, since $f(z)$ is not a polynomial.

Gelfond's proof was carried over to the case of a real quadratic irrationality $\beta$ by Kusmin in 1930. However, this method fails if $\beta$ is an algebraic irrationality of degree $h > 2$; it only gives the weaker result that at least one of the $h-1$ numbers $a^\beta, a^{\beta^2}, \dots, a^{\beta^{h-1}}$ is transcendental.

In 1934 Gelfond and Schneider, independently of each other, solved the general problem of the transcendency of $a^\beta$ for arbitrary irrational algebraic $\beta$. Both proofs make use of the arithmetical lemmas in \S 2, Chapter II, in order to construct suitable approximation forms. Since Schneider's proof is more closely related to the other ideas of Chapter II, we present it first. Gelfond's proof will be given in a simplified form which makes it somewhat shorter than Schneider's proof.

\section{Schneider's proof}\label{sec:1}

Suppose that $b$ is irrational and that the three numbers $b$, $a \neq 0$, $c = a^b = e^{b \log a}$ ($\log a \neq 0$) lie in an algebraic number field $K$ of degree $h$. If $a$ is a root of unity, then $c$ is not a root of unity; in this case, we replace $a, b, c$ by $c, b^{-1}, a$, so that the new $a$ is not a root of unity.

Define
\[
m = 4h + 3, \quad n = \frac{q^2}{m},
\]
where $q$ is a positive rational integer such that $q^2$ is divisible by $m$. We now apply \mylem{lem:2}, \S 2, Chapter II, to the determination of $m$ polynomials $P_1(x), \dots, P_m(x)$ of degree $\le 2n - 1$ with integral coefficients in $K$, not all $0$, such that the entire function
\[
R(x) = P_1(x)a^{(m-1)x} + P_2(x)a^{(m-2)x} + \dots + P_m(x)
\]
vanishes at the $q^2 = mn$ different points $x = \lambda + \mu b$ ($\lambda, \mu = 1, 2, \dots, q$). This condition yields $mn$ homogeneous linear equations for the $2mn$ indeterminate coefficients in $P_1, \dots, P_m$. The numerical coefficients of these equations lie in $K$ and the absolute values of all their conjugates are $\le (c_1 q)^{2n-1} (c_2^q)^m = O(c_3^n n^n)$, where $c_1, c_2, \dots$ denote positive rational integers independent of $n$; moreover, there exists a common denominator $c_4^{2n-1} c_5^{q(m-1)}$. In virtue of the lemma, we obtain a non-trivial solution such that all coefficients in $P_1, \dots, P_m$, and their conjugates, are $O(c_6^n n^n)$.

Putting
\begin{equation}
P_{kl}(x) = a^{(m-l)k} P_l(x+k) \qquad (k, l = 1, \dots, m), \tag{110}\label{eq:110}
\end{equation}
we have
\begin{equation}
R(x+k) = P_{k1}(x)a^{(m-1)x} + \dots + P_{km}(x). \tag{111}\label{eq:111}
\end{equation}

Suppose that $P_{\nu_1}(x), \dots, P_{\nu_g}(x)$ ($1 \le \nu_1 < \nu_2 < \dots < \nu_g \le m$) are not identically $0$, whereas all other $P_l(x) = 0$. If
\[
P_{\nu_l}(x) = \alpha_l x^{r_l} + \dots, \quad \alpha_l \neq 0 \quad (l=1,\dots,g),
\]
so that $r_l \ge 0$ denotes the exact degree of $P_{\nu_l}(x)$, then the $g$-rowed determinant
\[
\Delta(x) = | P_{k \nu_l}(x) |_{k,l=1,\dots,g} = \alpha_1 \dots \alpha_g v x^{r_1+\dots+r_g} + \dots,
\]
with
\[
v = |a^{(m-\nu_l)k}|_{k,l=1,\dots,g} = \prod_{k > l} (a^{m-\nu_l} - a^{m-\nu_k}) \neq 0,
\]
does not vanish identically, since $a$ is not a root of unity. The degree of $\Delta(x)$ is $r_1+\dots+r_g \le m(2n-1) < 2mn$. Therefore we can find at least one of the $2q^2 = 2mn$ numbers
\[
x = \lambda + \mu b \quad (\lambda = 1, \dots, 2q; \mu = 1, \dots, q),
\]
say $x = \xi$, such that
\[
\Delta(\xi) \neq 0.
\]
It follows from \eqref{eq:110} and \eqref{eq:111} that at least one of the $m$ numbers $R(\xi + k)$ ($k = 1, \dots, m$) is different from $0$. Suppose
\[
R(\xi + k) = \gamma \neq 0,
\]
then $c_4^{2n-1} c_5^{(m-1)(3q+m)} \gamma$ is an algebraic integer, and whence
\begin{equation}
|N(\gamma)| > c_7^{-n}. \tag{112}\label{eq:112}
\end{equation}

On the other hand, using our previous estimate of the coefficients of $P_l$, we obtain
\begin{equation}
\overline{|\gamma|} \le 2mn(c_8 q)^{2n-1} c_9^q O(c_6^n n^n) = O(c_{10}^n n^{2n}). \tag{113}\label{eq:113}
\end{equation}

It remains to find a sufficiently good upper estimate of $|\gamma|$ itself. Apply Cauchy's theorem to the entire function
\[
S(x) = R(x) \prod_{\lambda, \mu = 1}^{q} \frac{\xi+k - \lambda - \mu b}{x - \lambda - \mu b}
\]
which takes the value $\gamma$ at $x = \xi + k$; then
\[
\gamma = \frac{1}{2\pi i} \int_{C} \frac{S(x)}{x - \xi - k} \, dx,
\]
where $C$ is a simple closed curve, in positive direction, containing $\xi + k$ in its interior. Taking for $C$ the circle
\[
|x| = n + m + q(2+|b|) = n + O(\sqrt{n})
\]
we have
\[
|R(x)| \le 2mn(c_{11}n)^{2n-1} c_{12}^n O(c_6^n n^n) = O(c_{13}^n n^{3n}),
\]
\[
|x-\xi-k| \ge n > \frac{|x|}{c_{14}}, \qquad |x-\lambda-\mu b| > n \quad (\lambda, \mu=1,\dots,q),
\]
\[
\prod_{\lambda, \mu = 1}^{q} \frac{\xi+k-\lambda-\mu b}{x-\lambda-\mu b} = n^{-q^2} O((c_{15}q)^{q^2}) = O(c_{16}^n n^{-\frac{mn}{2}}),
\]
hence
\begin{equation}
\gamma = O(c_{13}^n n^{3n}) O(c_{16}^n n^{-\frac{mn}{2}}) = O(c_{17}^n n^{(3-\frac{m}{2})n}). \tag{114}\label{eq:114}
\end{equation}

In view of \eqref{eq:113} and \eqref{eq:114},
\begin{equation}
|N(\gamma)| = O\left(c_{18}^n n^{(3-\frac{m}{2})n+2(h-1)n}\right). \tag{115}\label{eq:115}
\end{equation}
Here the exponent of $n$ equals
\[
(3 - \frac{m}{2}) n + 2(h-1)n = -\frac{n}{2},
\]
so that \eqref{eq:112} and \eqref{eq:115} contradict each other as $n \to \infty$.

\section{Gelfond's proof}\label{sec:2}

We use $a, b, c, h, K$ in the same meaning as in \autoref{sec:1}: The three numbers $a, b$ and $c = a^b = e^{b \log a}$ lie in an algebraic number field of degree $h$; the number $b$ is irrational, $a$ and $\log a$ are $\neq 0$. However, we need not assume that $a$ is not a root of unity.

The approximation form is constructed in a different way. Put
\[
m = 2h + 2,
\]
\[
n = \frac{q^2}{2m},
\]
where $t = q^2$ is a positive rational integer square divisible by $2m$, and let
\[
\rho_1, \rho_2, \dots, \rho_t
\]
be the numbers
\[
(\lambda + \mu b) \log a \qquad (\lambda, \mu = 1, \dots, q).
\]
Introduce the entire function
\[
R(x) = \eta_1 e^{\rho_1 x} + \dots + \eta_t e^{\rho_t x}
\]
with indeterminate coefficients $\eta_1, \dots, \eta_t$ and consider the $mn$ homogeneous linear equations
\[
(\log a)^{-k} R^{(k)}(l) = 0 \qquad (k=0,\dots,n-1; l=1,\dots,m)
\]
for the $2mn = t$ unknown quantities $\eta_1, \dots, \eta_t$. Their numerical coefficients lie in $K$, they have a number
\[
c_1^{n-1+2mq} = O(c_2^n)
\]
as common denominator, and all their conjugates are $O((c_3 q)^{n-1} c_4^q) = O(c_5^n n^{\frac{n}{2}})$. It follows from \mylem{lem:2} that the equations have a non-trivial solution in integers $\eta_1, \dots, \eta_t$ of $K$ such that
\[
\overline{|\eta_k|} = O(c_6^n n^{\frac{n}{2}}) \qquad (k=1,\dots,t).
\]

Since the $t$ numbers $\rho_k$ all are different, the function $R(x)$ does not vanish identically. Choose $p$ such that $R^{(k)}(l) = 0$ for $k = 0, 1, \dots, p-1$ and $l = 1, \dots, m$, but $R^{(p)}(l) \neq 0$, for some $l = 1, \dots, m$. Clearly, $p \ge n$. Consider the number
\begin{equation}
(\log a)^{-p} R^{(p)}(l) = \gamma \neq 0; \tag{116}\label{eq:116}
\end{equation}
it lies in $K$, and $c_1^{p+2mq} \gamma$ is integral, so that
\begin{equation}
|N(\gamma)| > c_7^{-p}. \tag{117}\label{eq:117}
\end{equation}
Moreover,
\begin{equation}
\overline{|\gamma|} = t (c_8 q)^p c_9^q O(c_6^n n^{\frac{n}{2}}) = O(c_{10}^p p^{\frac{p}{2}}). \tag{118}\label{eq:118}
\end{equation}

To find again a suitable upper estimate of $|\gamma|$ itself, we introduce the entire function
\[
S(x) = p! \frac{R(x)}{(x-l)^p} \prod_{\substack{k=1\\k\neq l}}^m \left(\frac{l-k}{x-k}\right)^p;
\]
then
\[
\gamma = (\log a)^{-p} S(l)
\]
and
\[
S(l) = \frac{1}{2\pi i} \int_{C} \frac{S(x)}{x-l} \, dx,
\]
where we take for $C$ the circle
\[
|x| = m\left(1 + \frac{p}{q}\right).
\]
We obtain
\[
|R(x)| \le t c_{11}^{p+q} O(c_6^n n^{\frac{n}{2}}) = O(c_{12}^p p^{\frac{p}{2}}),
\]
\[
|x-l| \ge |x| \frac{p}{p+q}, \quad |x-k| \ge m \frac{p}{q} \quad (k=1,\dots,m),
\]
\[
|x-l|^{-p} \prod_{\substack{k=1 \\ k\neq l}}^{m} \left|\frac{l-k}{x-k}\right|^{p} = O\left(c_{13}^p p^{-p} \left(\frac{q}{p}\right)^{mp}\right),
\]
\[
|S(x)| \le p! c_{13}^p p^{-p} \left(\frac{q}{p}\right)^{mp} O\left(c_{12}^p p^{\frac{p}{2}}\right) = O\left(c_{14}^p p^{\frac{3-m}{2}p}\right);
\]
whence
\[
\gamma = O\left(c_{15}^p p^{\frac{3-m}{2}p}\right)
\]
and, by \eqref{eq:118},
\begin{equation}
|N(\gamma)| = O\left(c_{16}^p p^{(h-1)p + \frac{3-m}{2}p}\right). \tag{119}\label{eq:119}
\end{equation}
The exponent of $p$ equals
\[
(h-1)p + \frac{3-m}{2}p = -\frac{p}{2},
\]
so that \eqref{eq:117} and \eqref{eq:119} contradict each other for sufficiently large $n$, because of $p \ge n$.

\section{Additional remarks}\label{sec:3}

The main difference of the two proofs lies in the method by which the algebraic number $\gamma \neq 0$ is obtained. Schneider applies the functional equation $a^{x+y} = a^x a^y$ in order to construct $g$ approximation forms with non-vanishing determinant, and the inequality $\gamma \neq 0$ follows from an algebraical reason, in analogy to our procedure in \S 4 and \S 5, Chapter II. Gelfond simply uses the fact that the Taylor series of an analytic function which is not a polynomial, must contain infinitely many coefficients $\neq 0$. Since this method is indirect, it does not give immediately an explicit finite upper bound for $p$ in \eqref{eq:116}; though this can be found by a more detailed investigation. Gelfond's version of the proof is appropriate for the solution of the analogous problems concerning elliptic functions which we shall study in Chapter IV.

Our result on the transcendency of $a^b$ can also be stated this way: If $a$ and $c$ are algebraic numbers, $a, c \neq 0$, $\log a \neq 0$, then the ratio $\log c / \log a$ is either rational or transcendental. In other words: The logarithm of any algebraic number relative to any algebraic base is either rational or transcendental. However, it is not known whether $\log c$ and $\log a$ are algebraically independent in the case of an irrational quotient $\log c / \log a$; it is not even known whether there cannot exist an inhomogeneous linear relation $\alpha \log a + \gamma \log c = 1$ with quadratic irrational $\alpha$ and $\gamma$. Another example showing the narrow limits of our knowledge on transcendental numbers is the following one: Since $e$ is transcendental, not both numbers $e + \pi$ and $e \pi$ can be algebraic; but we do not even know whether $e + \pi$ or $e \pi$ are irrational.

Since $e^{\pi i} = -1$, the transcendency of $e^\pi$ is contained in Gelfond's result of 1929 on the transcendency of $a^b$ for algebraic $a \neq 0$, $\log a \neq 0$, and imaginary quadratic $b$. Before this discovery the problem of proving, for instance, the irrationality of $2^{\sqrt{2}}$ had been considered as extremely difficult, so that Hilbert liked to mention it as a problem whose solution lay still further in the future than the proof of Riemann's hypothesis or Fermat's conjecture. This shows that one cannot guess the real difficulties of a problem before having solved it.
